
% --- Comment out chunks ---
\usepackage{comment}

% --- Core Accessibility & Structure ---
\usepackage[american]{babel}
\tagpdfsetup{math/alt/use,activate,interwordspace=true,uncompress}


% --- Math Fonts and others ---
\usepackage{graphicx}
\usepackage{booktabs}
\usepackage{tabularx}
\usepackage{xcolor}
\usepackage{fancyvrb}
\usepackage{tikz}
\usetikzlibrary{calc, shapes, backgrounds}
\frenchspacing

% --- Hyperlinks ---
\usepackage{hyperref}
\hypersetup{
	pdftitle={Lecture Template},
	colorlinks=true,%
	citecolor=blue,%
	filecolor=blue,%
	linkcolor=blue,%
	urlcolor=blue
}

% --- Color Definitions (Defined Once) ---
\definecolor{UFBlue}{HTML}{0021A5}
\definecolor{UFOrange}{HTML}{A53200} 
\definecolor{UFGreen}{HTML}{1A5632}
% ======================================================================
% FRAME DESIGN & COLOR CONFIGURATION
% ======================================================================
% 1. Style the container (The Banner Background)
\EditInstance{header}{std}{
	background-color = UFBlue,
	color = white,
	print-frame-title = true
}

% 2. Style the text inside the container (The Typography)
\EditInstance{frametitle}{header}{
	font = \bfseries\huge
}

% 3. Banner Visibility Toggles 
% a. Turn banner off (Make it transparent and heightless)
\newcommand{\banneroff}{%
	\EditInstance{header}{std}{
		background-color = white, % Matches slide background
		color = white,            % Hides title text
		print-frame-title = false
	}
}
% b. Restore standard Blue Banner
\newcommand{\banneron}{%
	\EditInstance{header}{std}{
		background-color = UFBlue,
		color = white,
		print-frame-title = true
	}
}

% Frame Numbering
\makeatletter
% 1. Create a custom metadata variable containing the formatted string.
%    Placing \hfill inside the variable forces it to the right,
%    overriding the footer's default alignment.
\newcommand{\@slideprogress}{\hfill \@framenumber/\@totalframes}
\makeatother

% 2. Register the new variable as a valid footer element in ltx-talk.
\DeclareInstance{footer-element}{slideprogress}{talk}{}

% 3. Apply the single element to the footer configuration.
\EditInstance{footer}{std}{
	background-color = white,
	color = black,
	element-order = {slideprogress},
	left-hspace = 5mm,
	right-hspace = 5mm,
	font = \small\sffamily
}

% =====================================================================
% MULTI-SLIDE EXAMPLE ENVIRONMENT (ORANGE BANNER)
% ======================================================================
% --- UF BRANDING TOGGLES ---
% Use these to switch banner colors for Examples or Special Sections
\newcommand{\ufblue}{%
	\EditInstance{header}{std}{background-color = UFBlue}%
}
\newcommand{\uforange}{%
	\EditInstance{header}{std}{background-color = UFOrange}%
}

% Alias for pedagogical clarity
\let\exampleon\uforange
\let\exampleoff\ufblue

% ======================================================================
% CUSTOM SEMANTIC BLOCKS (Corrected Left Border)
% =====================================================================
\newenvironment{theory}[1][]{%
	\vspace{0.5em}\noindent
	\textcolor{UFBlue}{\vrule width 3pt}\hspace{0.5em}%
	\begin{minipage}[t]{\dimexpr\textwidth-3pt-0.5em\relax}
		\begin{block}{\strut\color{UFBlue}\bfseries Theory\if\relax\detokenize{#1}\relax\else: #1\fi}
			\vspace{2pt}}{%
			\vspace{2pt}\end{block}
	\end{minipage}\vspace{0.5em}}  

\newenvironment{example}[1][]{%
	\vspace{0.5em}\noindent
	\textcolor{UFOrange}{\vrule width 3pt}\hspace{0.5em}%
	\begin{minipage}[t]{\dimexpr\textwidth-3pt-0.5em\relax}
		\begin{block}{\strut\color{UFOrange}\bfseries Example\if\relax\detokenize{#1}\relax\else: #1\fi}
			\vspace{2pt}}{%
			\vspace{2pt}\end{block}
	\end{minipage}\vspace{0.5em}}  

\newenvironment{remark}[1][]{%
	\vspace{0.5em}\noindent
	\textcolor{UFGreen}{\vrule width 3pt}\hspace{0.5em}%
	\begin{minipage}[t]{\dimexpr\textwidth-3pt-0.5em\relax}
		\begin{block}{\strut\color{UFGreen}\bfseries Remark\if\relax\detokenize{#1}\relax\else: #1\fi}
			\vspace{2pt}}{%
			\vspace{2pt}\end{block}
	\end{minipage}\vspace{0.5em}}  

% ======================================================================
% GLOBAL HOOKS & OVERLAYS
% ======================================================================
\makeatletter
\renewcommand{\@listi}{%
	\leftmargin\leftmargini
	\parsep 0pt
	\topsep 0pt
	\itemsep 8pt % Increased for better readability
}
\let\@listI\@listi
\makeatother

\NewDocumentCommand\fade{D<>{all}}{%
	\temporal<#1>{\color{black!40}}{\color{black}}{\color{black!10}}
}

% Automated Section Transitions
\let\oldsection\section
\renewcommand{\section}[1]{%
	\oldsection{#1}
	\begin{frame}
		\frametitle{Outline}
		\tableofcontents[currentsection]
	\end{frame}
}